\section{The coupled co-design problem}
\label{sec:codesign}

\subsection{Statement}

The algebraic layer of \cref{sec:closure} assumed the whole mission is spent
hovering. The differential layer of \cref{sec:dynamics} assumed the mass was
known. Neither is true independently. The honest statement couples them.

\begin{problem}[Co-design]
\label{prob:codesign}
Find a design vector $p$ and a control law $u(\cdot)$ such that
\begin{align}
  \dt{x} &= f(x,u,w;p), \qquad x(0)=x_0(p), \label{eq:cd1}\\
  m &= \mfix + \mstr(p) + \mprop(p) + \mbat, \label{eq:cd2}\\
  \eb(1-\varsigma)\,\mbat &\ge \int_0^{t_f} P\big(x(t),u(t);p\big)\dd t ,
  \qquad
  p_b\,\mbat \ge \max_{t\in[0,t_f]} P\big(x(t),u(t);p\big), \label{eq:cd3}
\end{align}
where in this section $\eb$ denotes the usable specific energy before the
landing reserve $\varsigma$ is removed and $p_b$ the pack's specific power,
subject to the path constraints
\begin{equation}
  T(t)\le\lambda mg,\quad
  \Omega_i(t)\le\Omega_{\max},\quad
  \ell(t)\ge d_{\mathrm{safe}},\quad
  \theta_{\mathrm{read}}(t)\le\theta_{\max},
  \label{eq:cd4}
\end{equation}
with $\ell(t)$ the distance from the user's eye to the nearest rotor disk,
minimising a scalar objective
\begin{equation}
  \mathcal{J}(p,u) = w_m m + w_P \bar P + w_N L_p + w_s\,\mathrm{span}
                     + w_e\, e_{\mathrm{lag}} .
  \label{eq:cd5}
\end{equation}
\end{problem}

\Cref{eq:cd3} is the coupling. The right-hand side depends on the trajectory,
which depends on the plant, which depends on $m$, which contains $\mbat$, which
is the left-hand side. The first inequality keeps the reserve: starting from
$E(0)=\eb\mbat$, the state of charge at $t_f$ is at least $\varsigma$. For a
minimally sized continuous battery one of the two inequalities is an equality,
\begin{equation}
  \mbat = \max\left\{\frac{1}{\eb(1-\varsigma)}\int_0^{t_f}P\,\dd t,\;
                     \frac{1}{p_b}\max_t P\right\} ;
  \label{eq:batreq}
\end{equation}
a discrete pack, built from whole cells, need not satisfy either with equality.

\subsection{Reduction to a fixed point}

\Cref{eq:cd2,eq:cd3} can be reduced to a scalar fixed-point problem, which is
what makes the system tractable.

\begin{definition}[Inner and outer maps]
For a given battery mass $\mu$, let
\begin{equation}
  \mathcal{G}(\mu) \coloneqq \text{the solution } m \text{ of }
  m = \mfix + \mstr(m) + \mprop(m) + \mu
  \label{eq:innermap}
\end{equation}
be the gross mass consistent with that battery, and let
\begin{equation}
  \mathcal{H}(\mu) \coloneqq
  \max\left\{
  \frac{1}{\eb(1-\varsigma)}\int_0^{t_f}
    P\big(x(t;\mathcal{G}(\mu),\mu),u(t)\big)\dd t,\;
  \frac{1}{p_b}\max_t P\big(x(t;\mathcal{G}(\mu),\mu),u(t)\big)\right\}
  \label{eq:outermap}
\end{equation}
be the battery mass that the resulting mission demands, by \cref{eq:batreq}.
\end{definition}

\begin{proposition}[Well-posedness of the inner map]
\label{prop:inner}
Let $S(m)\coloneqq\mstr(m)+\mprop(m)$ map $[0,\infty)$ into $[0,\infty)$ (masses
are non-negative) and satisfy the Lipschitz bound
\begin{equation}
  |S(m_2)-S(m_1)| \le L\,|m_2-m_1|
  \quad\text{for all } m_1,m_2\ge0,
  \qquad L<1 .
  \label{eq:innercond}
\end{equation}
Then for every $\mu\ge0$, \cref{eq:innermap} has exactly one solution
$\mathcal{G}(\mu)$, the iteration $m\leftarrow\mfix+S(m)+\mu$ converges to it
from any start at rate $L$, and $\mathcal{G}$ is Lipschitz with constant
$1/(1-L)$. If $S$ is non-decreasing, so is $\mathcal{G}$. If $S$ is
differentiable,
\begin{equation}
  \mathcal{G}'(\mu) = \frac{1}{1-S'\big(\mathcal{G}(\mu)\big)} .
  \label{eq:Gprime}
\end{equation}
\end{proposition}

\begin{proof}
The map $\Phi(m)=\mfix+S(m)+\mu$ sends $[0,\infty)$ into itself because
$S\ge0$, and is a contraction with constant $L$ by \cref{eq:innercond}, so the
Banach fixed-point
theorem~\cite[Ch.~5]{ortega1970} gives existence, uniqueness and convergence of
the iteration. For $\mu_1,\mu_2$ with solutions $m_1,m_2$,
$m_2-m_1 = S(m_2)-S(m_1)+\mu_2-\mu_1$, so
$|m_2-m_1|\le L|m_2-m_1|+|\mu_2-\mu_1|$, which is the Lipschitz bound. If $S$ is
non-decreasing and $\mu_2>\mu_1$ but $m_2<m_1$, then
$m_2-m_1-\big(S(m_2)-S(m_1)\big)\le-(1-L)(m_1-m_2)<0$, contradicting
$m_2-m_1-\big(S(m_2)-S(m_1)\big)=\mu_2-\mu_1>0$. Finally
$\Xi(m,\mu)=\Phi(m)-m$ has $\partial_m\Xi = S'-1\le L-1<0$, so the implicit
function theorem applies and gives \cref{eq:Gprime}.
\end{proof}

\begin{remark}[Why a slope bound and not a growth bound]
Continuity, monotonicity and $\limsup S(m)/m<1$ give existence but not
uniqueness. The smooth, increasing, bounded
$S(m)=4/[1+e^{-4(m-3)}]$ with $\mfix+\mu=1$ makes $m=1+S(m)$ hold at
$m\approx1.0013$, $3$ and $4.9987$. The slope condition is what excludes this.
For the lumped models $S'=f_s+\gamma_m$. For the component models of
\cref{sec:components}, evaluated on the design of \cref{sec:results}, $S'$ lies
between $0.13$ and $0.23$ over \SIrange{0.8}{50}{\kilo\gram}; the test suite
asserts $0\le S'<0.9$ on that interval. That is a check on one design over a
bounded interval, not a proof for every parameter set.
\end{remark}

\begin{theorem}[Existence of a closed design]
\label{thm:codesignfixed}
The system \cref{eq:cd2,eq:cd3} is equivalent to the scalar fixed-point equation
\begin{equation}
  \mu = \mathcal{H}(\mu) .
  \label{eq:fixedpoint}
\end{equation}
If $\mathcal{H}$ is continuous and maps a compact interval $[0,\mu_{\max}]$ into
itself, a solution exists by Brouwer's theorem. If in addition $\mathcal{H}$ is
a contraction, $\mathrm{Lip}(\mathcal{H}) = L < 1$, the solution is unique and
the iteration $\mu_{k+1}=\mathcal{H}(\mu_k)$ converges geometrically with rate
$L$ from any starting point.
\end{theorem}

\begin{proof}
Equivalence is by construction. The existence and uniqueness statements are
Brouwer's fixed-point theorem and the Banach fixed-point theorem
respectively~\cite[Ch.~5]{ortega1970}.
\end{proof}

\begin{proposition}[Contraction criterion]
\label{prop:contraction}
Suppose the mission power is well approximated by a mission-dependent multiple
of the hover power, $\bar P = \varkappa\,P_{\mathrm{hover}}(m)$ with $\varkappa$
independent of $m$, and that the energy term attains the maximum in
\cref{eq:outermap} in a neighbourhood of $\mu$. Then
\begin{equation}
  \mathcal{H}'(\mu)
  = \frac{\varkappa\,t_f}{\eb(1-\varsigma)}\,
    \frac{\dd P_{\mathrm{hover}}}{\dd m}\,\mathcal{G}'(\mu),
  \label{eq:Hprime}
\end{equation}
and in the fixed-rotor case with the lumped models, where
$\mathcal{G}'=1/(1-f_s-\gamma_m)=1/\alpha$,
\begin{equation}
  \mathcal{H}'(\mu) = \frac{\varkappa}{\alpha}\cdot
    \frac{3}{2}\,\beta\,\sqrt{m} ,
  \label{eq:Hprime2}
\end{equation}
with $\beta$ as in \cref{eq:beta}, whose specific energy is net of the reserve,
$\eb(1-\varsigma)$ in the notation of this section. The iteration therefore
contracts precisely when $\tfrac32\varkappa\beta\sqrt m < \alpha$, which at
$\varkappa=1$ is the condition $F'(m)>0$ of \cref{cor:dmdm}, i.e.\ $m<m^*$. On
the power-limited branch $\mathcal{H}' = p_b^{-1}(\dd P_{\mathrm{peak}}/\dd m)
\mathcal{G}'$ instead, and where the two branches exchange the maximum
$\mathcal{H}$ has a kink.
\end{proposition}

\begin{proof}
Differentiate the energy term of \cref{eq:outermap} by the chain rule. From
\cref{eq:beta,eq:ptot},
$P_{\mathrm{hover}}(m) = \beta\,\eb(1-\varsigma)\, m^{3/2}/t_f + \Paux$, so
$\dd P_{\mathrm{hover}}/\dd m = \tfrac32\beta\,\eb(1-\varsigma)\, m^{1/2}/t_f$.
Substituting into \cref{eq:Hprime} with $\mathcal{G}'=1/\alpha$, the factors
$t_f$ and $\eb(1-\varsigma)$ cancel and \cref{eq:Hprime2} remains.
\end{proof}

\begin{remark}[What the overhead does, and does not, do]
\Cref{prop:contraction} connects the algebraic and the simulated pictures: the
condition for the simulation-in-the-loop iteration to converge is the
light-branch condition with $\beta$ multiplied by $\varkappa$. It is tempting to
say that overhead ``degrades $\alpha$''. That is a reading of one inequality,
not an identity. Multiplying the whole power by $\varkappa$ turns the closure
into
\begin{equation}
  \alpha m - \varkappa\beta\, m^{3/2}
  = \mfix + \varkappa\,\frac{\Paux t_f}{\eb},
  \label{eq:kappaclosure}
\end{equation}
which scales the energy coupling and the auxiliary part of $\Mzero$ together,
and by \cref{eq:closedform} lowers the critical load by the factor
$1/\varkappa^2$. A change of $\alpha$ would do neither.
\end{remark}

\subsection{Residual formulation}

Rather than iterate \cref{eq:fixedpoint} naively, the engine treats the pair
$y = (m,\mbat)$ and drives the residual
\begin{equation}
  \vect{\mathcal{R}}(y) =
  \begin{bmatrix}
    m - \big(\mfix + \mstr(m) + \mprop(m) + \mbat\big) \\[2pt]
    \mbat - \max\left\{\dfrac{E_{\mathrm{mission}}(m,\mbat)}{\eb(1-\varsigma)},\;
                       \dfrac{P_{\mathrm{peak}}(m,\mbat)}{p_b}\right\}
  \end{bmatrix}
  = \vect 0 .
  \label{eq:residual}
\end{equation}
The first component is algebraic and cheap; the second requires a trajectory
simulation. The engine solves $\mathcal{R}_1=0$ at every outer step by the
iteration of \cref{prop:inner}, reducing the problem to the scalar
$\mathcal{R}_2=0$ at one simulation per update. When $|\mathcal{R}_2|$ falls
below tolerance the battery is rounded up to the last requirement and the
design is flown again; the run is declared converged only when that
verification flight requires no more battery than it carries. The energy margin
$\eb(1-\varsigma)\mbat-E_{\mathrm{mission}}$, the power margin
$p_b\mbat-P_{\mathrm{peak}}$ and the path constraints \cref{eq:cd4} of the final
flight are then reported separately. Closing the mass and energy balance and
satisfying the path constraints are different outcomes: a converged battery
loop can still be an unacceptable flight.

\subsection{The outer optimisation}

With the closure solved, \cref{eq:cd5} becomes a function of the design vector
alone, and \cref{prob:codesign} reduces to a constrained finite-dimensional
program
\begin{equation}
  \min_{p\in\mathcal{P}} \ \mathcal{J}(p)
  \quad\text{s.t.}\quad
  g_j(p)\le0,\ j=1,\dots,J ,
  \label{eq:outerprog}
\end{equation}
with $\mathcal{P}$ a box and $g_j$ the constraints of \cref{eq:cd4} together
with the modelling-validity constraints of
\cref{as:sigma,as:stall,eq:reynolds}. The objective is non-convex, cheap to
evaluate on the algebraic layer and expensive on the simulated one, and has
large infeasible regions where the closure has no solution at all. This
combination dictates the choice of a population-based method with exterior
penalties, discussed in \cref{sec:numerics}.

\begin{remark}[Why not adjoint optimal control]
\Cref{prob:codesign} in full generality is a control co-design problem and
could be attacked with adjoint-based methods over the trajectory. We do not,
for two reasons. First, the objective is dominated by steady hover, for which
the trajectory contributes a scalar $\varkappa$ near unity; the sensitivity of
the design to the control law is small compared with its sensitivity to the
rotor geometry. Second, and more importantly, the constraints that actually
bind, acoustic and geometric, are algebraic functions of $p$ and do not involve
$u(\cdot)$ at all. Spending effort on trajectory optimisation would refine the
part of the problem that is not deciding the answer.
\end{remark}
